\section*{अध्याय \ref{ch:g11:polarity-and-cohesion} --- विद्युत-ऋणात्मकता, ध्रुवता और अंतराआण्विक बल}

\begin{solution}{exo:g11:polarity-and-cohesion:1}
\ce{H-Cl}: क्लोरीन (2.20 के मुक़ाबले 3.16)। \ce{O-H}: ऑक्सीजन (3.44)।
\ce{C-O}: ऑक्सीजन (2.55 के मुक़ाबले 3.44)।
\end{solution}

\begin{solution}{exo:g11:polarity-and-cohesion:2}
\ce{H-H}: 0, ध्रुवीय नहीं। \ce{C-H}: 0.35, अध्रुवीय माना जाता है। \ce{O-H}:
1.24, ध्रुवीय। \ce{N-H}: 0.84, ध्रुवीय। \ce{Cl-Cl}: 0, ध्रुवीय नहीं। \ce{C-Cl}:
0.61, ध्रुवीय।
\end{solution}

\begin{solution}{exo:g11:polarity-and-cohesion:3}
नाइट्रोजन (3.04 > 2.19) और ऑक्सीजन (3.44 > 2.58): \omterm{def:g11:polarity-and-cohesion:electronegativity}{विद्युत-ऋणात्मकता}
समूह में नीचे घटती है। नाइट्रोजन (3.04 > 2.55) और मैग्नीशियम
(1.31 > 0.93): वह \omterm{def:g9:periodic-table-first-look:period}{आवर्त} में बाएँ से दाएँ बढ़ती है।
\end{solution}

\begin{solution}{exo:g11:polarity-and-cohesion:4}
अमोनिया और पानी: हर एक में नाइट्रोजन या ऑक्सीजन से आबंधित हाइड्रोजन \omterm{def:g7:atoms-and-molecules:atom}{परमाणु} हैं,
और उसी \omterm{def:g7:atoms-and-molecules:atom}{परमाणु} पर \omterm{def:g10:lewis-and-shape:lone-pair}{एकाकी युग्म}। मेथेन में N, O या F पर कोई हाइड्रोजन नहीं है;
हाइड्रोजन क्लोराइड में हाइड्रोजन क्लोरीन से आबंधित है, जो परिभाषा के तीन \omterm{def:g7:atoms-and-molecules:atom}{परमाणुओं} में
से नहीं है।
\end{solution}

\begin{solution}{exo:g11:polarity-and-cohesion:5}
हाँ: \omterm{def:g11:polarity-and-cohesion:van-der-waals}{वान्डरवाल्स अन्योन्यक्रियाएँ} सभी \omterm{def:g7:atoms-and-molecules:molecule}{अणुओं} के बीच होती हैं। पर मेथेन के
\omterm{def:g7:atoms-and-molecules:molecule}{अणु} छोटे हैं (10 \omterm{def:g9:inside-the-atom:nucleus}{इलेक्ट्रॉन}), इसलिए ये अन्योन्यक्रियाएँ बहुत दुर्बल हैं
और बहुत कम ताप पर टूट जाती हैं: मेथेन लगभग
\qty{-162}{\celsius} पर उबलती है।
\end{solution}

\begin{solution}{exo:g11:polarity-and-cohesion:6}
\ce{CH2Cl2}: ध्रुवीय (दो ध्रुवीय \ce{C-Cl} आबंध और दो लगभग अध्रुवीय
\ce{C-H} आबंध: कोई सममिति उन्हें निरस्त नहीं करती)। \ce{CCl4}: अध्रुवीय (चतुष्फलक
के कोनों पर चार एक जैसे आबंध)। \ce{NH3}: ध्रुवीय
(पिरामिड)। \ce{CO2}: अध्रुवीय (रैखिक, विपरीत आबंध)। \ce{BF3}:
अध्रुवीय (एक तल में \ang{120} पर तीन एक जैसे \omterm{def:g11:polarity-and-cohesion:polar-bond}{ध्रुवीय आबंध} निरस्त हो जाते हैं)।
\end{solution}

\begin{solution}{exo:g11:polarity-and-cohesion:7}
\ce{O-H} समूह का हाइड्रोजन दिया जा सकता है (\ce{CH3} के हाइड्रोजन नहीं,
जो कार्बन से आबंधित हैं); अपने दो \omterm{def:g10:lewis-and-shape:lone-pair}{एकाकी युग्मों} वाला ऑक्सीजन \omterm{def:g7:atoms-and-molecules:atom}{परमाणु}
ग्रहण कर सकता है। हाँ: मेथेनॉल का \ce{O-H} पानी के \omterm{def:g7:atoms-and-molecules:molecule}{अणु} के ऑक्सीजन से आबंध बना
सकता है, और पानी का एक \ce{O-H} मेथेनॉल के ऑक्सीजन से, उदाहरण के लिए
\ce{CH3-O-H}\,$\cdots$\,\ce{OH2}।
\end{solution}

\begin{solution}{exo:g11:polarity-and-cohesion:8}
तीनों \omterm{def:g7:atoms-and-molecules:molecule}{अणु} अध्रुवीय हैं और उनकी आकृति एक ही है, पर वे उत्तरोत्तर
बड़े हैं (10, 18 और 36 \omterm{def:g9:inside-the-atom:nucleus}{इलेक्ट्रॉन}): उनकी \omterm{def:g11:polarity-and-cohesion:van-der-waals}{वान्डरवाल्स अन्योन्यक्रियाएँ} बढ़ती हैं,
और उनके साथ उनके क्वथन ताप भी।
\end{solution}

\begin{solution}{exo:g11:polarity-and-cohesion:9}
\omterm{def:g9:periodic-table-first-look:period}{समूह} 14: मेथेन में N, O या F से आबंधित कोई हाइड्रोजन नहीं है और वह कोई
\omterm{def:g11:polarity-and-cohesion:hydrogen-bond}{हाइड्रोजन आबंध} नहीं बनाती, इसलिए वह अपने समूह की प्रवृत्ति का पालन करती है।
\end{solution}

\begin{solution}{exo:g11:polarity-and-cohesion:10}
\omterm{def:g11:polarity-and-cohesion:van-der-waals}{वान्डरवाल्स अन्योन्यक्रियाएँ}: \ce{HCl} से \ce{HI} तक \omterm{def:g7:atoms-and-molecules:molecule}{अणु} बड़े होते जाते
हैं, और आकार के साथ इन अन्योन्यक्रियाओं की वृद्धि \omterm{def:g11:polarity-and-cohesion:polar-bond}{आंशिक आवेशों} के बीच के
आकर्षण की कमी पर भारी पड़ती है।
\end{solution}

\begin{solution}{exo:g11:polarity-and-cohesion:11}
आयोडीन का \omterm{def:g7:atoms-and-molecules:molecule}{अणु} क्लोरीन के \omterm{def:g7:atoms-and-molecules:molecule}{अणु} से कहीं बड़ा है (34 के मुक़ाबले 106
\omterm{def:g9:inside-the-atom:nucleus}{इलेक्ट्रॉन}): उसकी \omterm{def:g11:polarity-and-cohesion:van-der-waals}{वान्डरवाल्स अन्योन्यक्रियाएँ} कहीं प्रबल हैं,
इतनी प्रबल कि कमरे के ताप पर \omterm{def:g7:atoms-and-molecules:molecule}{अणुओं} को ठोस में थामे रखें।
\end{solution}

\begin{solution}{exo:g11:polarity-and-cohesion:12}
प्रोपेन अध्रुवीय है: केवल \omterm{def:g11:polarity-and-cohesion:van-der-waals}{वान्डरवाल्स अन्योन्यक्रियाएँ}। मेथॉक्सीमेथेन
ध्रुवीय है (दो ध्रुवीय \ce{C-O} आबंध, एक मुड़े हुए \ce{C-O-C} में), पर उसके ऑक्सीजन पर
कोई हाइड्रोजन नहीं है: उसके \omterm{def:g7:atoms-and-molecules:molecule}{अणुओं} के बीच कोई \omterm{def:g11:polarity-and-cohesion:hydrogen-bond}{हाइड्रोजन आबंध} नहीं। एथेनॉल
ध्रुवीय है और उसके \ce{O-H} समूह \omterm{def:g11:polarity-and-cohesion:hydrogen-bond}{हाइड्रोजन आबंध} बनाते हैं। आकर्षण जितने प्रबल,
क्वथन ताप उतना ऊँचा: प्रोपेन < मेथॉक्सीमेथेन
< एथेनॉल।
\end{solution}

\begin{solution}{exo:g11:polarity-and-cohesion:13}
बर्फ़ का खुला जाल उन्हीं \omterm{def:g7:atoms-and-molecules:molecule}{अणुओं} की तुलना में अधिक जगह घेरता है जितनी वे द्रव में
घेरते हैं, इसलिए बर्फ़ द्रव पानी से कम घनी है और तैरती है। झील ऊपर से जमती
है: बर्फ़ की परत नीचे के पानी को ऊष्मारोधन देती है, जो
द्रव बना रहता है, और मछलियाँ जीवित रहती हैं।
\end{solution}

\begin{solution}{exo:g11:polarity-and-cohesion:14}
ठोस मेथेन (छोटे \omterm{def:g7:atoms-and-molecules:molecule}{अणुओं} के बीच \omterm{def:g11:polarity-and-cohesion:van-der-waals}{वान्डरवाल्स अन्योन्यक्रियाएँ}) <
बर्फ़ (\omterm{def:g11:polarity-and-cohesion:hydrogen-bond}{हाइड्रोजन आबंध}) < पोटैशियम क्लोराइड (\omterm{def:g9:ions:ion}{आयनों} के बीच आकर्षण)।
\end{solution}

\begin{solution}{exo:g11:polarity-and-cohesion:15}
पानी के \omterm{def:g7:atoms-and-molecules:molecule}{अणु} के पास देने के लिए दो हाइड्रोजन \omterm{def:g7:atoms-and-molecules:atom}{परमाणु} और ग्रहण करने के लिए दो \omterm{def:g10:lewis-and-shape:lone-pair}{एकाकी
युग्म} हैं: हर \omterm{def:g7:atoms-and-molecules:molecule}{अणु} चार \omterm{def:g11:polarity-and-cohesion:hydrogen-bond}{हाइड्रोजन आबंधों} में भाग ले सकता है, दो
दिए हुए और दो ग्रहण किए हुए, जिससे एक पूरा जाल बनता है। हाइड्रोजन फ़्लुओराइड के
\omterm{def:g7:atoms-and-molecules:molecule}{अणु} के पास तीन \omterm{def:g10:lewis-and-shape:lone-pair}{एकाकी युग्म} हैं पर केवल एक हाइड्रोजन \omterm{def:g7:atoms-and-molecules:atom}{परमाणु}: वह
एक ही \omterm{def:g11:polarity-and-cohesion:hydrogen-bond}{हाइड्रोजन आबंध} देता है, इसलिए औसतन प्रति \omterm{def:g7:atoms-and-molecules:molecule}{अणु} केवल एक आबंध
बन सकता है (टेढ़ी-मेढ़ी श्रृंखलाएँ)। पानी में प्रति \omterm{def:g7:atoms-and-molecules:molecule}{अणु} लगभग दुगने \omterm{def:g11:polarity-and-cohesion:hydrogen-bond}{हाइड्रोजन आबंध} हैं,
और वह अधिक ऊँचे ताप पर उबलता है।
\end{solution}

\begin{solution}{pb:g11:polarity-and-cohesion:1}
\textbf{1.} \ce{H-O-H}, \ce{H-S-H}, \ce{H-Se-H}, हर केंद्रीय \omterm{def:g7:atoms-and-molecules:atom}{परमाणु} पर
दो \omterm{def:g10:lewis-and-shape:lone-pair}{एकाकी युग्म}: O, S और Se में छह \omterm{def:g10:electron-shells:valence}{संयोजकता इलेक्ट्रॉन} हैं, क्योंकि वे एक ही
समूह में हैं।

\textbf{2.} $M(\ce{H2O}) = \qty{18.0}{g/mol}$, $M(\ce{H2S}) =
\qty{34.1}{g/mol}$, $M(\ce{H2Se}) = \qty{81.0}{g/mol}$।

\textbf{3.} $3.44 - 2.20 = 1.24$; $2.58 - 2.20 = 0.38$;
$2.55 - 2.20 = 0.35$। केवल \ce{O-H} आबंध स्पष्ट रूप से ध्रुवीय है।

\textbf{4.} मुड़ी हुई (चार समूह, जिनमें दो \omterm{def:g10:lewis-and-shape:lone-pair}{एकाकी युग्म})। पानी ध्रुवीय है;
हाइड्रोजन सल्फ़ाइड और हाइड्रोजन सेलेनाइड बहुत थोड़े ही।

\textbf{5.} 10, 18 और 36 \omterm{def:g9:inside-the-atom:nucleus}{इलेक्ट्रॉन}: हाइड्रोजन सेलेनाइड की
\omterm{def:g11:polarity-and-cohesion:van-der-waals}{वान्डरवाल्स अन्योन्यक्रियाएँ} सबसे प्रबल हैं।

\textbf{6.} केवल पानी: उसके हाइड्रोजन \omterm{def:g7:atoms-and-molecules:atom}{परमाणु} ऑक्सीजन से आबंधित हैं, जो तीन बहुत
विद्युत-ऋणात्मक \omterm{def:g7:atoms-and-molecules:atom}{परमाणुओं} में से एक है; सल्फ़र और सेलेनियम इतने
विद्युत-ऋणात्मक नहीं कि अपने हाइड्रोजन \omterm{def:g7:atoms-and-molecules:atom}{परमाणुओं} को स्पष्ट
$\delta^+$ दे सकें।

\textbf{7.} चार: दो अपने हाइड्रोजन \omterm{def:g7:atoms-and-molecules:atom}{परमाणुओं} से, दो अपने
दो \omterm{def:g10:lewis-and-shape:lone-pair}{एकाकी युग्मों} से ग्रहण किए हुए।

\textbf{8.} \omterm{def:g11:polarity-and-cohesion:van-der-waals}{वान्डरवाल्स अन्योन्यक्रियाएँ} (छोटे \omterm{def:g11:polarity-and-cohesion:polar-bond}{आंशिक आवेशों} के बीच के
दुर्बल आकर्षण सहित)।

\textbf{9.} $212.9 - 273.15 = \qty{-60.3}{\celsius}$ और
\qty{-41.3}{\celsius}।

\textbf{10.} $-41.3 - (-60.3) = \qty{19.0}{\celsius}$।

\textbf{11.} \omterm{def:g7:atoms-and-molecules:molecule}{अणु} बड़ा है, इसलिए उसकी \omterm{def:g11:polarity-and-cohesion:van-der-waals}{वान्डरवाल्स अन्योन्यक्रियाएँ}
अधिक प्रबल हैं।

\textbf{12.} $-60.3 - 19.0 = \qty{-79.3}{\celsius}$।

\textbf{13.} वृद्धि $-88.1 - (-112) = \qty{23.9}{\celsius}$; मेथेन के लिए
अनुमान $-112 - 23.9 = \qty{-136}{\celsius}$, जबकि वास्तव में \qty{-162}{\celsius}:
बहिर्वेशन लगभग \qty{26}{\celsius} अधिक है। यह
सटीक नहीं है; वास्तविक प्रवृत्ति सीधी रेखा नहीं है।

\textbf{14.} वृद्धि $\qty{18.7}{\celsius}$; हाइड्रोजन फ़्लुओराइड के लिए अनुमान
$-85.1 - 18.7 = \qty{-103.8}{\celsius}$, जबकि वास्तव में \qty{19.6}{\celsius}: लगभग
\qty{123}{\celsius} का अंतर।

\textbf{15.} $100 - (-79.3) = \qty{179}{\celsius}$, लगभग
\qty{180}{\celsius}।

\textbf{16.} वृद्धि $\qty{25.3}{\celsius}$; अनुमान $-87.8 - 25.3 =
\qty{-113}{\celsius}$; अंतर $-33.4 - (-113) = \qty{80}{\celsius}$।

\textbf{17.} पानी (\qty{180}{\celsius}) > हाइड्रोजन फ़्लुओराइड
(\qty{123}{\celsius}) > अमोनिया (\qty{80}{\celsius})। पानी दो
हाइड्रोजन \omterm{def:g7:atoms-and-molecules:atom}{परमाणु} देता है और दो \omterm{def:g10:lewis-and-shape:lone-pair}{एकाकी युग्मों} से ग्रहण करता है: प्रति \omterm{def:g7:atoms-and-molecules:molecule}{अणु} चार
\omterm{def:g11:polarity-and-cohesion:hydrogen-bond}{हाइड्रोजन आबंध}, सभी उपयोग में। हाइड्रोजन फ़्लुओराइड के पास देने के लिए एक हाइड्रोजन
है, और अमोनिया के पास ग्रहण करने के लिए केवल एक \omterm{def:g10:lewis-and-shape:lone-pair}{एकाकी युग्म}: औसतन प्रति \omterm{def:g7:atoms-and-molecules:molecule}{अणु} कम
आबंध, और नाइट्रोजन, जो कम विद्युत-ऋणात्मक है, दुर्बल आबंध बनाता है।

\textbf{18.} गैस: वह लगभग \qty{-80}{\celsius} पर उबलता। पृथ्वी पर
कोई द्रव पानी न होता, न महासागर, न बारिश, और न वह जीवन
जो उन पर निर्भर है।

\textbf{19.} यह केवल दो बिंदुओं से खींची गई सीधी रेखा का बहिर्वेशन करता है,
और \omterm{def:g9:periodic-table-first-look:period}{समूह} 14 पर परीक्षण दिखाता है कि ऐसा बहिर्वेशन
लगभग \qty{25}{\celsius} तक चूक सकता है।

\textbf{20.} \omterm{def:g11:polarity-and-cohesion:hydrogen-bond}{हाइड्रोजन आबंधों} के बिना पानी लगभग \qty{-80}{\celsius} पर उबलता,
अपने वास्तविक क्वथनांक से लगभग \qty{180}{\celsius}
नीचे।
\end{solution}
